%%%%%%%%%%  scan idx 239  |  folio 226  %%%%%%%%%%
\SGAPrintPage{226}
We may therefore suppose that \(Y=Z[T_1,\ldots,T_r]\), with \(f\) the canonical projection. The diagram \((*)\) then corresponds to a diagram of rings
\[
\begin{tikzcd}
& B=C[T_1,\ldots,T_r] \arrow[dl,"u"'] & \\
A & C \arrow[l,"v"] \arrow[u,"\varepsilon"'] & {,}
\end{tikzcd}
\]
where \(u\) and \(v\) are epimorphisms, and \(\varepsilon\) the canonical inclusion. For \(1\leq i\leq r\), put \(u(T_i)=x_i\), and let \(z_i\in C\) satisfy \(v(z_i)=x_i\). The retraction \(\sigma:B\longrightarrow C\) of \(\varepsilon\) defined by \(\sigma(T_i)=z_i\) satisfies \(v\sigma=u\). In geometric language, it corresponds to a section \(s\) of \(f\) such that \(i=sj\). By (VII 1.10), \(s\) is a regular immersion, so the assertion follows from (1.1.1).

\uline{Proof of} (1.1)(\uline{end}). Let \(f:X\to Y\) be a locally finite type morphism. There exists an open covering \((X_\alpha)\) of \(X\) such that each \(X_\alpha\) embeds by a closed immersion into a scheme \(Y_\alpha\) smooth over \(Y\). Putting \(\coprod X_\alpha=X'\), \(\coprod Y_\alpha=Y'\), we thus have a commutative diagram
\begin{equation}\tag{$a'$}
\begin{tikzcd}
X' \arrow[r,"i'"] \arrow[d,"p'"'] & Y' \arrow[ddl,"f'"]\\
X \arrow[d,"f"']\\
Y
\end{tikzcd}
\end{equation}
where \(p'\) is a morphism that is the sum of a surjective family of open immersions (we shall say \uline{special} for short), \(i'\) a closed immersion, and \(f'\) a smooth morphism. We shall define \(D(f)_{n\text{-}\coh}\) as the full subcategory of \(D(X)\) formed by the objects \(E\) such that \(i'_*p'^*(E)\) is \(n\)-pseudo-coherent. Let us first show that this definition makes sense, i.e. that \(D(f)_{n\text{-}\coh}\) is independent of the choice of \(i'\), \(p'\), \(f'\). Thus, let

%%%%%%%%%%  scan idx 240  |  folio 227  %%%%%%%%%%
\SGAPrintPage{227}
\begin{equation}\tag{$a''$}
\begin{tikzcd}
X'' \arrow[r,"i''"] \arrow[d,"p''"'] & Y'' \arrow[ddl,"f''"]\\
X \arrow[d,"f"']\\
Y
\end{tikzcd}
\end{equation}
be a commutative diagram analogous to \((a')\) (\(p''\) a special morphism, \(i''\) a closed immersion, \(f''\) a smooth morphism). From \((a')\) and \((a'')\), we deduce a commutative diagram
\[
\begin{tikzcd}
X'\times_X X'' \arrow[rr,hook] \arrow[dr] \arrow[dd] & & Y'\times_Y Y'' \arrow[dr] \arrow[dd] & \\
& X' \arrow[rr,hook,"i'"] \arrow[dd] & & Y' \arrow[dd,"f'"]\\
X'' \arrow[rr,hook,"i''"] \arrow[dr] & & Y'' \arrow[dr,"f''"'] & \\
& X \arrow[rr,"f"] & & Y,
\end{tikzcd}
\]
where the horizontal arrows, except for \(f\), are closed immersions, and the arrows of the left (resp. right) Cartesian square are special morphisms (resp. smooth morphisms). We are therefore reduced to proving the following assertion:

(A): \uline{Let}
\[
\begin{tikzcd}
X_1 \arrow[r,hook,"i_1"] \arrow[d,"p"'] & Y_1 \arrow[d,"q"]\\
X_2 \arrow[r,hook,"i_2"] & Y_2
\end{tikzcd}
\]
\uline{be a commutative square of schemes}, \uline{where} \(p\) \uline{is a special morphism}, \(q\) \uline{a smooth morphism}, \(i_1\) \uline{and} \(i_2\) \uline{closed immersions}, \uline{and let} \(E\in\operatorname{ob} D(X_2)\). \uline{Then}, \uline{for} \(i_{2*}(E)\) \uline{to be} \(n\)-\uline{pseudo-coherent}, \uline{it is necessary and sufficient that} \(i_{1*}p^*(E)\) \uline{be so}.

%%%%%%%%%%  scan idx 241  |  folio 228  %%%%%%%%%%
\SGAPrintPage{228}
\uline{Proof of} (A). By definition, \(p\) is the sum of a family of open subsets \(U_\alpha\) of \(X_2\) that cover \(X_2\). Each \(U_\alpha\) is induced by an open subset \(V_\alpha\) of \(Y_2\), and the \(V_\alpha\) cover an open neighborhood \(X'_2\) of \(X_2\) in \(Y_2\). We therefore have a commutative diagram with a Cartesian square
\[
\begin{tikzcd}
X_1 \arrow[r,hook,"i'_1"] \arrow[d,"p"'] & X'_1 \arrow[d,"p'"]\\
X_2 \arrow[r,hook,"i'_2"] \arrow[rr,bend right,"i_2"'] & X'_2 \arrow[r,hook,"j"] & Y_2,
\end{tikzcd}
\]
where \(i'_1\) and \(i'_2\) are closed immersions, \(p\) and \(p'\) special morphisms, and \(j\) an open immersion. The condition ``\(i_{2*}(E)\) \(n\)-pseudo-coherent'' is clearly equivalent to ``\(i'_{2*}(E)\) \(n\)-pseudo-coherent''. On the other hand, it follows at once from (1.1.3) that the condition ``\(i'_{2*}(E)\) \(n\)-pseudo-coherent'' is equivalent to ``\(i'_{1*}p^*(E)\) \(n\)-pseudo-coherent''. Applying (1.1.4) to the factorizations \((jp')i'_1\) and \(qi_1\) of \(i_2p\), and to \(p^*(E)\in\Ob D(X_1)\), we finally conclude that ``\(i_{2*}(E)\) \(n\)-pseudo-coherent'' is equivalent to ``\(i_{1*}p^*(E)\) \(n\)-pseudo-coherent'', q.e.d.

We have thus shown that the category \(D(f)_{n\text{-}\coh}\) defined above by means of the diagram \((a')\) in fact depends only on \(f\). Moreover, it is trivial from the definition that it does satisfy conditions (i) to (iii) of (1.1). This therefore completes the proof of (1.1).

\sgaStableTarget{sga6.e03.def.1-2}\textbf{Definition 1.2.} \uline{Let} \(f:X\to Y\) \uline{be a locally finite type morphism}, \(E\in\Ob D(X)\), \(n\in\mathbb Z\). \uline{We shall say that} \(E\) \uline{is} \(n\)-pseudo-coherent relative to \(f\) \uline{if} \(E\in\Ob D(f)_{n\text{-}\coh}\) (cf. (1.1)). \uline{We shall say that} \(E\) \uline{is} pseudo-coherent (\uline{or equivalently} \((-\infty)\)-pseudo-coherent) relative to \(f\) \uline{if} \(E\) \uline{is} \(n\)-\uline{pseudo-coherent relative to} \(f\) \uline{for every} \(n\in\mathbb Z\). \uline{We shall write} \(D(f)_{\mathrm{coh}}\) \uline{for the full subcategory}

%%%%%%%%%%  scan idx 242  |  folio 229  %%%%%%%%%%
\SGAPrintPage{229}
\uline{of} \(D(X)\) \uline{formed by the objects pseudo-coherent relative to} \(f\). \uline{We shall say that} \(f\) \uline{is} \(n\)-\uline{pseudo-coherent} (\uline{resp.} \uline{pseudo-coherent}) \uline{if} \(\underline{O}_X\) \uline{is} \(n\)-\uline{pseudo-coherent} (\uline{resp.} \uline{pseudo-coherent}) \uline{relative to} \(f\). \uline{We shall sometimes say}, \uline{when this cannot lead to confusion}, ``\(n\)-\uline{pseudo-coherent} (\uline{resp.} \uline{pseudo-coherent}) \uline{relative to} \(Y\)'' \uline{instead of} ``\uline{relative to} \(f\)'', \uline{and shall also write} \(D(X)_{Y\text{-}n\text{-}\coh}\) (\uline{resp.} \(D(X)_{Y\text{-}\coh}\)) \uline{instead of} \(D(f)_{n\text{-}\coh}\) (\uline{resp.} \(D(f)_{\mathrm{coh}}\)). \uline{Finally, for} \(*=+,-,b\), \uline{we shall put}
\[
\begin{gathered}
  D^*(f)_{n\text{-}\coh}=D^*(X)\cap D(f)_{n\text{-}\coh},\\
  D^*(f)_{\mathrm{coh}}=D^*(X)\cap D(f)_{\mathrm{coh}}.
\end{gathered}
\]

\sgaStableTarget{sga6.e03.rem.1-2-1}\textbf{Remarks 1.2.1.}
\begin{enumerate}[label=\alph*)]
\item The categories \(D(f)_{n\text{-}\coh}\), \(D(f)_{\mathrm{coh}}\) are not in general contained in \(D^{-}(X)\). This is so, however, when \(X\) is quasi-compact, since an object of \(D(X)\) with locally bounded-above cohomology then has bounded-above cohomology.

\item The objects of \(D(f)_{\mathrm{coh}}\) have quasi-coherent cohomology. Indeed, let \(E\in\Ob D(f)_{\mathrm{coh}}\); the question being local, we may suppose that \(f\) is a closed immersion. Then \(f_*(E)\) is pseudo-coherent, hence has quasi-coherent cohomology, and the result follows from the fact that an \(\underline{O}_X\)-Module \(M\) is quasi-coherent if and only if \(f_*(M)\) is so.

\item When speaking of an \(\underline{O}_X\)-Module, we shall use the terminology ``of finite presentation rel. to \(f\)'' in preference to ``\((-1)\)-pseudo-coherent rel. to \(f\)''. For an \(\underline{O}_X\)-Module \(E\) to be of finite presentation rel. to \(f\), it is necessary and sufficient, in view of (I 2.9), that there exist locally a closed immersion \(i:X\to X'\), with \(X'\) smooth over \(Y\), such that \(i_*(E)\) is of finite presentation (in the ordinary sense). In particular, saying that \(f\) is

%%%%%%%%%%  scan idx 243  |  folio 230  %%%%%%%%%%
\SGAPrintPage{230}
\((-1)\)-pseudo-coherent is equivalent to saying that \(f\) is locally of finite presentation in the sense of (EGA IV 1.4).
\end{enumerate}

\sgaStableTarget{sga6.e03.ex.1-3}\textbf{Examples 1.3.} A smooth morphism and a regular immersion (1.1.2) are pseudo-coherent morphisms. The same holds for any locally finite type morphism \(f:X\to Y\) when \(Y\) is locally noetherian.

\sgaStableTarget{sga6.e03.prop.1-4}\textbf{Proposition 1.4.} \uline{Let} \(f:X\to Y\) \uline{be a locally finite type morphism}, \uline{and let}
\[
\begin{tikzcd}
& G \arrow[dl] & \\
E \arrow[rr] && F \arrow[ul]
\end{tikzcd}
\]
\uline{be a distinguished triangle in} \(D(X)\). \uline{We then have the following table}, \uline{where} \(n\in\mathbb Z\cup\{-\infty\}\) \uline{indicates a degree of pseudo-coherence relative to} \(f\), \uline{and each row represents an implication whose conclusion is underlined}:
\[
\begin{array}{ccc}
E&F&G\\
n+1&n&\underline n\\
n&\underline n&n\\
\underline n&n&n-1.
\end{array}
\]
\uline{In particular}, \(D(f)_{\mathrm{coh}}\) \uline{is a triangulated subcategory of} \(D(X)\).

\uline{Proof}. By properties (i) to (iii) of (1.1), we may suppose that \(f\) is a closed immersion. We then have a distinguished triangle
\[
\begin{tikzcd}
& f_*(G) \arrow[dl] & \\
f_*(E) \arrow[rr] && f_*(F) \arrow[ul]
\end{tikzcd},
\]
and the assertion follows from (I 2.6).

\sgaStableTarget{sga6.e03.def.1-5}\textbf{Definition 1.5.} \uline{Let} \(X\) \uline{and} \(Y\) \uline{be two} \(S\)-\uline{schemes}. \uline{We shall say that} \(X\) \uline{and} \(Y\) \uline{are tor-independent over} \(S\) \uline{if we have}
\[
  \underline{\Tor}_i^S(\underline{O}_X,\underline{O}_Y)=0\quad\text{\uline{for} }i\ne 0,
\]
\uline{where the} \(\underline{\Tor}_i^S\) \uline{are the local hyper-tor defined in} (EGA III 6.5).

